% =====================================================================
%  PAPER III -- from the CERTIFLIGHT concept (MSCA O3/T3/T4, WP4).
%  Target: NeurIPS (main track).
%
%  Standalone. Set as Overleaf's main document; pdfLaTeX + BibTeX.
%    pdflatex PaperIII_certified_constrained_marl && bibtex PaperIII_certified_constrained_marl \
%      && pdflatex PaperIII_certified_constrained_marl && pdflatex PaperIII_certified_constrained_marl
% =====================================================================
\documentclass[11pt,a4paper]{article}
% =====================================================================
%  Shared preamble for the three research proposals.
%
%  Each proposal \input{}s this file and is otherwise standalone, so any
%  one of them can be set as Overleaf's main document and compiled on its
%  own. Nothing here is specific to a single proposal.
% =====================================================================
\usepackage[T1]{fontenc}
\usepackage[utf8]{inputenc}
% Times, matching the parent manuscript. mathptmx is in every TeX Live and on
% Overleaf; lmodern is deliberately not used, since it is absent from minimal
% distributions and this preamble should compile anywhere.
\usepackage{mathptmx}
\usepackage[a4paper,margin=25mm]{geometry}
\usepackage{amsmath,amssymb,amsthm}
\usepackage{booktabs}
\usepackage{tabularx}
\usepackage{xcolor}
\usepackage{titlesec}
\usepackage{enumitem}
\usepackage{microtype}
\usepackage[hidelinks,breaklinks]{hyperref}

% ---- shared style tokens, matching the parent manuscript ------------
\definecolor{netcol}{RGB}{31,78,121}
\definecolor{autocol}{RGB}{18,120,90}
\definecolor{bridgecol}{RGB}{176,96,0}
\definecolor{accent}{RGB}{140,20,20}

\newcommand{\layernet}{\textcolor{netcol}{\textbf{network}}}
\newcommand{\layerauto}{\textcolor{autocol}{\textbf{autonomy}}}
\newcommand{\dataset}[1]{\textsc{#1}}
\newcommand{\code}[1]{\texttt{#1}}
\newcommand{\Real}{\mathbb{R}}
\newcommand{\Norm}[1]{\left\lVert #1\right\rVert}
\newcommand{\MCR}{\mathrm{MCR}}

\newtheorem{claim}{Claim}
\newtheorem{proposition}{Proposition}
\newtheorem{definition}{Definition}
\theoremstyle{remark}
\newtheorem{remark}{Remark}

% ---- section styling -------------------------------------------------
\titleformat{\section}{\large\bfseries\color{netcol}}{\thesection}{0.7em}{}
\titleformat{\subsection}{\normalsize\bfseries}{\thesubsection}{0.6em}{}
\titlespacing*{\section}{0pt}{1.4em}{0.6em}

\setlist[itemize]{leftmargin=1.35em,topsep=0.3em,itemsep=0.25em}
\setlist[enumerate]{leftmargin=1.6em,topsep=0.3em,itemsep=0.25em}

% ---- a boxed "what this rests on" callout ---------------------------
\newenvironment{honestly}
  {\par\medskip\noindent\begin{tabular}{@{}p{2.5mm}@{\hspace{3mm}}p{0.9\linewidth}@{}}
   \textcolor{accent}{\rule{2.5mm}{\baselineskip}} &
   \small\itshape}
  {\end{tabular}\par\medskip}

% ---- the proposal header --------------------------------------------
\newcommand{\proposalheader}[3]{%
  \begin{center}
  {\Large\bfseries #1\par}\medskip
  {\large\itshape #2\par}\medskip
  {\small #3\par}
  \end{center}
  \medskip\hrule\medskip}

\usepackage{array}
\newcolumntype{Y}{>{\raggedright\arraybackslash}X}
\newcolumntype{P}[1]{>{\raggedright\arraybackslash}p{#1}}
\newcommand{\Rmin}{R_{\min}}
\newcommand{\Ri}{R_i}
\newcommand{\prov}[1]{\textsc{\small[#1]}}
% Quantifier-level labels (defined in full in Paper II; referenced here).
\newcommand{\Qlvl}[1]{\textbf{Q}\textsubscript{\textit{#1}}}
\newcommand{\Qop}{\Qlvl{op}}
\newcommand{\Qpop}{\Qlvl{pop}}

\begin{document}

\proposalheader
  {Paper III --- Certified-Constrained MARL}
  {Certificates as Constraints: Feasibility-Preserving Multi-Agent Reinforcement
   Learning for Swarms Whose Guarantees Degrade}
  {Research proposal from the CERTIFLIGHT concept (WP4) $\cdot$ RobustUAVs.ai
   $\cdot$ Roger Nick Anaedevha $\cdot$ 2026-10-01\\
   Target: NeurIPS, main track $\cdot$ builds on the parent
   manuscript~\cite{parentpaper} and Papers~I--II}

\begin{abstract}
\noindent
Safe reinforcement learning offers two kinds of guarantee and neither is the one
a safety-critical fleet needs. Lagrangian methods satisfy constraints
\emph{in expectation} over a horizon; barrier and shielding methods give
\emph{pointwise} guarantees but read their constraint sets from hand-specified
margins. We study the case where the constraint set is instead the output of a
formal robustness certificate that is itself \emph{time-varying and
observation-dependent}: each agent carries a certified navigation radius
$\Ri(t)$ that shrinks as the agent is jammed or faulted, and may hold a task only
while $\Ri(t)$ exceeds the radius that task provably requires. This couples a
certificate to a constrained Markov decision process in a way the safe-RL
literature has not analysed, because the feasible set moves as a function of the
observation the policy is still learning to act on. We give three results: a
projection safety layer that preserves pointwise feasibility at every step of
\emph{both learning and deployment} whenever a feasible assignment exists; an
exact, checkable condition for when one exists, via a deficiency form of Hall's
theorem; and a bound on coverage loss as a function of the fraction of agents
whose certificate has collapsed, matching a brute-force oracle in simulation. The
contribution is not a new RL algorithm but a correctness analysis of
certificate-as-constraint, with the certificate taken as an adversarially
moving object rather than a fixed box.
\end{abstract}

% =====================================================================
\section{The problem the existing guarantees miss}
% =====================================================================
A swarm flies a coverage mission --- urban-corridor monitoring, environmental
transects, point-to-point medical delivery. Each aircraft's navigation stack
carries a certificate (the parent manuscript's~\cite{parentpaper}, or any
producing a radius $\Ri(t)$: the largest position error under which the aircraft
still completes its assigned task). Under jamming or a sensor fault, $\Ri(t)$
shrinks. Below the radius a task requires, the aircraft can no longer be trusted
with it, and the fleet must re-allocate.

The decision --- which aircraft holds which task, as certificates move --- is a
constrained sequential control problem. The safe-RL literature has the pieces
and not the composition:

\begin{itemize}
\item \textbf{Lagrangian constrained RL}~\cite{achiam2017cpo,paternain2019zeroduality}
  satisfies constraints \emph{in expectation over the horizon}. A fleet that
  violates a safety radius on $5\%$ of steps but satisfies it on average is not
  acceptable; expectation is the wrong quantifier for flight.
\item \textbf{Shielding and control barrier functions}~\cite{alshiekh2018shielding,ames2019cbf}
  give pointwise guarantees, but from a \emph{hand-specified} safe set. The
  margin is an engineering choice disconnected from any formal statement about
  the perception stack.
\item \textbf{Certified RL}~\cite{wu2022crop,yang2024carol} certifies a policy
  against an observation-perturbation ball, single-agent, with the ball fixed in
  advance. It never feeds a certificate into a \emph{fleet-level} constraint.
\end{itemize}

\noindent
The gap is specific: \textbf{no analysis treats the constraint set of a
multi-agent policy as the output of a formal certificate that varies with the
observation.} That is the object here, and it is harder than a fixed constraint
because the feasible region the policy must stay inside is itself a function of
the jamming the policy is reacting to.

% =====================================================================
\section{Setup}
% =====================================================================
A fleet of $N$ agents and $M$ tasks is a constrained Markov decision
process~\cite{altman1999cmdp} $(\mathcal S,\mathcal A,P,r,\{c_i\})$. Agent $i$'s
state includes its certified radius $\Ri(t)$, computed online by the perception
certificate from locally estimated jamming and fault conditions. An allocation
$\pi$ assigns agents to tasks to maximise coverage, subject to the pointwise
constraints
\begin{equation}
\label{eq:constraint}
c_i:\quad \Ri(t)\;\ge\;\Rmin(\tau_i)\quad\text{whenever agent $i$ holds task
$\tau_i$, at every $t$.}
\end{equation}
We take $\Ri(t)$ as given by the certificate and treat it, adversarially, as an
arbitrary bounded cadlag process --- no distributional assumption, because a
jammer is not a distribution. $\Rmin(\tau)$ is the task's required radius, from
the mission geometry.

% =====================================================================
\section{Three results}
% =====================================================================
\subsection{Feasibility preservation (T3)}
\begin{proposition}[Projection preserves pointwise feasibility]
\label{pr:feas}
Let $\mathcal F(t)$ be the set of assignments satisfying
\eqref{eq:constraint} at the current radii $\Ri(t)$. If the policy's proposed
assignment $a_t$ is replaced by its projection $\Pi_{\mathcal F(t)}(a_t)$ onto
the feasible set whenever $\mathcal F(t)\neq\varnothing$, then every executed
assignment satisfies
\eqref{eq:constraint} pointwise, during learning and deployment alike, regardless
of the policy's parameters.
\end{proposition}
\noindent
The content is not the projection --- that is immediate once $\mathcal F(t)$ is
known --- but that it holds \emph{during learning}, so a half-trained policy
never flies an infeasible assignment, and that $\mathcal F(t)$ is cheap to
represent (next section). This is what lets a certificate replace a hand-tuned
barrier margin: the margin becomes a computed, auditable quantity.

\subsection{When a feasible assignment exists (T3, exactly)}
Feasibility is a bipartite matching: an edge joins agent $i$ to task $j$ iff
$\Ri(t)\ge\Rmin(\tau_j)$. A full assignment exists iff the matching saturates the
tasks, and Hall's theorem~\cite{hall1935} gives the condition. We use its
deficiency form, specialised to the \emph{nested} structure that
\eqref{eq:constraint} forces --- an agent that qualifies for a demanding task
qualifies for every easier one --- which collapses the exponential Hall check to
a linear one.

\begin{proposition}[Checkable feasibility]
\label{pr:hall}
Sort task requirements descending, $\rho_1\ge\cdots\ge\rho_M$, and let
$n(\rho)=|\{i:\Ri(t)\ge\rho\}|$. The maximum number of tasks simultaneously
coverable is $M-d$, where the deficiency is
\[
d \;=\; \max_{1\le k\le M}\ \big(k-n(\rho_k)\big)^{+}.
\]
A full assignment exists iff $d=0$, checkable in $O(N+M\log M)$.
\end{proposition}

\begin{proposition}[$k$-robust coverage]
\label{pr:krobust}
Full coverage survives the collapse of any $k$ agents' certificates iff
$n(\rho_j)\ge j+k$ for every $j$. The largest such $k$ is the fleet's coverage
margin: the number of simultaneous certificate failures it tolerates with no
coverage loss.
\end{proposition}

\begin{honestly}
Both propositions are elementary once the nested structure is seen; the
contribution is seeing it, and that it makes an otherwise NP-flavoured
fleet-robustness question a linear-time check an onboard planner can run every
control step. We verified both against a brute-force matching oracle over 3000
random instances: exact agreement on the coverage count and on the $k$-robust
condition. That is a sanity check on the math, not evidence about flight.
\end{honestly}

\subsection{Coverage loss as certificates degrade (T4)}
\begin{proposition}[Coverage-degradation bound]
\label{pr:coverage}
If a fraction $\phi$ of agents have $\Ri(t)$ below the median task requirement,
coverage loss is at most $\lceil\phi N\rceil$ tasks, and exactly
$M-(M-d_\phi)$ with $d_\phi$ the deficiency of Proposition~\ref{pr:hall}
recomputed on the degraded radii. The resulting curve, coverage versus $\phi$,
is the operational trade-off between per-agent robustness investment and the
fleet size needed for a target coverage guarantee.
\end{proposition}
\noindent
Classical coverage control has results for \emph{failed} (binary) agents
\cite{cortes2004coverage}; the novelty is a bound parameterised by
\emph{continuous} certified-capability degradation, which is what a procurement
decision turns on --- how many aircraft, carrying how good a navigation stack, to
hold a coverage SLA under a given jamming level.

% =====================================================================
\section{What is proposed, and what is already in hand}
% =====================================================================
\begin{center}\small
\begin{tabularx}{\linewidth}{@{}P{0.34\linewidth}Y@{}}
\toprule
\textbf{In hand} & \textbf{To build} \\
\midrule
A certificate producing $\Ri(t)$ from jamming/fault state (parent manuscript;
$\gamma=1.625$\,m/s campaign supremum) &
the CMDP, the projection layer, and the learned allocation policy \\
A physics-informed fleet simulator, \code{MAX\_FLEET}=12, ten attack classes &
a jamming/fault/gust schedule driving $\Ri(t)$ per agent over the mission \\
Propositions \ref{pr:feas}--\ref{pr:coverage}, brute-force-checked &
the empirical study: feasibility violations, coverage curves, scaling \\
\bottomrule
\end{tabularx}
\end{center}

% =====================================================================
\section{Experimental plan}
% =====================================================================
\begin{description}[leftmargin=2.6em,style=nextline]
\item[RQ1] Does the projection layer achieve \emph{zero} pointwise violations
  during training, where a PPO-Lagrangian baseline~\cite{yu2022mappo,achiam2017cpo}
  achieves only in-expectation satisfaction? Report per-step violation counts,
  not episode averages.
\item[RQ2] How does the measured coverage-versus-$\phi$ curve compare to the
  bound of Proposition~\ref{pr:coverage}, across jamming intensities and fault
  and gust profiles?
\item[RQ3] How does the certified-constraint policy compare to
  jamming-blind and worst-case-fixed-margin baselines on coverage and on
  completion?
\item[RQ4] At what fleet size and task count does the projection step stop being
  real-time? Report the scalability ceiling rather than hiding it.
\end{description}

\paragraph{Baselines.} Jamming-blind MAPPO~\cite{yu2022mappo}; PPO-Lagrangian
and CPO~\cite{achiam2017cpo}; shielded MADDPG~\cite{lowe2017maddpg,alshiekh2018shielding};
QMIX~\cite{rashid2018qmix}; and a worst-case fixed-margin allocator. The point of
the comparison is to separate what the \emph{learned policy} achieves from what
the \emph{guarantee mechanism} enforces --- so every method is reported both with
and without the projection layer.

\paragraph{Why NeurIPS and not a robotics venue.} The result is about the
quantifier on a learned constraint --- pointwise versus in-expectation, under a
constraint set that is itself a certified, adversarially moving object. That is a
learning-theory and safe-RL contribution. The UAV swarm is the setting that makes
it concrete and gives it a benchmark, not the contribution itself.

% =====================================================================
\section{Relation to the other papers and to the parent}
% =====================================================================
\begin{itemize}
\item \textbf{Parent manuscript / Paper~I} supply $\Ri(t)$. Paper~III consumes it.
  The interface is exactly the certified radius; nothing about how it is computed
  leaks into the CMDP analysis, which is why Paper~III stands if the certificate
  is swapped.
\item \textbf{Paper~II (SoK)} places this work at its \Qop{}/\Qpop{} boundary:
  pointwise feasibility is a per-operation guarantee, the coverage bound a
  per-population one. Paper~III is a worked instance of a sound \Qop$\to$\Qpop{}
  lift.
\item \textbf{Distinct from all three} in object: Papers I, II and the parent are
  about a certificate; Paper~III is about a \emph{controller that consumes}
  one.
\end{itemize}

% =====================================================================
\section{Risks}
% =====================================================================
\begin{itemize}
\item \textbf{The analysis could outrun the experiments.} The three propositions
  are the contribution and are already proved and checked; the RL study is
  supporting evidence. If MARL training is unstable at the target fleet size, the
  paper reports the projection layer's guarantees (which hold for \emph{any}
  policy, including a random one) and the scalability ceiling, and still stands.
\item \textbf{$\Ri(t)$ is only as good as the certificate.} Paper~III assumes it;
  if the certificate is loose, constraints are conservative and coverage
  suffers, but feasibility and the bounds are unaffected. This dependence is
  stated, not hidden.
\item \textbf{Sim-to-real.} The study is in simulation. The host's UAV testbed is
  the path to a hardware vignette; it is a strengthener, not a dependency, and
  over-claiming transfer from simulation is the failure mode to avoid.
\item \textbf{Novelty challenge (``this is just matching + shielding'').} Pre-empt
  it: the novelty is the \emph{coupling} --- a formal perception certificate as a
  time-varying CMDP constraint with feasibility preserved \emph{during learning}
  --- and the linear-time feasibility and $k$-robustness conditions that the
  nested structure of certified radii makes possible.
\end{itemize}

% =====================================================================
\section{Page plan and timeline}
% =====================================================================
\begin{center}\small
\begin{tabular}{@{}lr@{\qquad}lr@{}}
\toprule
\textbf{Section} & \textbf{pp} & \textbf{Section} & \textbf{pp} \\
\midrule
Introduction                 & 1.0  & Experiments                 & 2.5 \\
Related work                 & 1.0  & Discussion and limitations  & 0.75 \\
Setup (CMDP, certificate)    & 1.0  & Conclusion                  & 0.25 \\
Theory (Prop.\ 1--4)         & 2.0  & (appendix: proofs, uncapped)& --- \\
\midrule
\multicolumn{3}{@{}l}{\textbf{Body, within a 9-page NeurIPS limit}} & \textbf{8.5} \\
\bottomrule
\end{tabular}
\end{center}

\noindent
About four to five months, and the theory is done, so writing can begin while the
RL study runs. The simulator and certificate exist; the new engineering is the
CMDP wrapper, the projection layer and the degradation schedule.

\small
\bibliographystyle{plain}
\bibliography{proposals}

\end{document}
